% ============================================================
% Section 83: 一维晶体中空穴的物理量
% ============================================================
\section{固体理论：一维晶体中空穴的物理量}
\label{sec:hole-properties}

\begin{modelbox}[题目：满带中移去电子后空穴各物理量的计算]
某一维晶体满带中电子能量关系为
\[
E(k)=\frac{\hbar^2k_1^2}{6m}-\frac{3\hbar^2k^2}{m},
\]
其中 $m$ 为电子质量，$k_1=\dfrac{1}{8a}$，晶格常数 $a=3.14\times10^{-10}\,\text{m}$。若从 $k_1$ 处的态中移去一个电子，其他所有态均被占据，$k_1$ 处出现一个空穴。试计算：
\begin{enumerate}
    \item 空穴的波矢 $k_h$；
    \item 空穴的有效质量 $m_h^*$；
    \item 空穴的速度 $v_h$；
    \item 空穴的动量；
    \item 空穴的能量 $E_h$；
    \item 若满带有相当部分电子被激发走，画出空穴带并标出 $k_1$ 处的空穴。
\end{enumerate}
\end{modelbox}

\begin{tipbox}[考点快速分析]
\begin{itemize}
    \item \textbf{空穴波矢}：$k_h=-k_e$（总波矢守恒，满带总波矢为零）
    \item \textbf{空穴有效质量}：$m_h^*=-m_e^*>0$（带顶电子 $m_e^*<0$）
    \item \textbf{空穴速度}：$v_h=v_e$（缺失电子的速度）
    \item \textbf{空穴能量}：$E_h=-E_e$（以带顶为零点，向下增加）
    \item \textbf{空穴动量}：$p_h=\hbar k_h=m_h^*v_h$（两种等效表示）
\end{itemize}
\end{tipbox}

\begin{derivationbox}[标准解题过程]
\textbf{预备计算}
\begin{equation}
k_1=\frac{1}{8a}=\frac{1}{8\times3.14\times10^{-10}}\approx3.98\times10^8\,\text{m}^{-1}.
\end{equation}

\textbf{(1) 空穴的波矢}

满带总波矢为零。移去波矢为 $k_e=k_1$ 的电子后，剩余电子总波矢为 $-k_1$。为保持总波矢守恒，空穴波矢
\begin{equation}
\boxed{k_h=-k_e=-k_1\approx-3.98\times10^8\,\text{m}^{-1}.}
\end{equation}
方向与 $k_1$ 相反。

\textbf{(2) 空穴的有效质量}

电子有效质量
\begin{equation}
m_e^*=\frac{\hbar^2}{d^2E/dk^2}=\frac{\hbar^2}{-6\hbar^2/m}=-\frac{m}{6}.
\end{equation}

空穴有效质量定义为电子有效质量的相反数：
\begin{equation}
\boxed{m_h^*=-m_e^*=\frac{m}{6}\approx\frac{9.11\times10^{-31}}{6}\approx1.52\times10^{-31}\,\text{kg}\approx0.17\,m_0.}
\end{equation}
为正值。

\textbf{(3) 空穴的速度}

电子在 $k_1$ 处的群速度
\begin{equation}
v_e=\frac{1}{\hbar}\frac{dE}{dk}\bigg|_{k=k_1}=\frac{1}{\hbar}\Bigl(-\frac{6\hbar^2k_1}{m}\Bigr)=-\frac{6\hbar k_1}{m}
\approx-2.77\times10^5\,\text{m/s}.
\end{equation}

空穴速度等于缺失电子的速度：
\begin{equation}
\boxed{v_h=v_e\approx-2.77\times10^5\,\text{m/s}.}
\end{equation}

\textbf{(4) 空穴的动量}

由波矢定义：
\begin{equation}
p_h=\hbar k_h=1.055\times10^{-34}\times(-3.98\times10^8)\approx-4.20\times10^{-26}\,\text{N}\cdot\text{s}.
\end{equation}

由有效质量和速度验证：
\begin{equation}
p_h=m_h^*v_h=1.52\times10^{-31}\times(-2.77\times10^5)\approx-4.20\times10^{-26}\,\text{N}\cdot\text{s}.
\end{equation}
两者一致。
\begin{equation}
\boxed{p_h\approx-4.20\times10^{-26}\,\text{N}\cdot\text{s}.}
\end{equation}

\textbf{(5) 空穴的能量}

能带顶在 $k=0$ 处，电子能量
\begin{equation}
E(0)=\frac{\hbar^2k_1^2}{6m}\approx3.20\times10^{-22}\,\text{J}\approx2.0\times10^{-3}\,\text{eV}.
\end{equation}

在 $k_1$ 处电子能量
\begin{equation}
E(k_1)=\frac{\hbar^2k_1^2}{6m}-\frac{3\hbar^2k_1^2}{m}=-\frac{17\hbar^2k_1^2}{6m}.
\end{equation}

空穴能量定义为缺失电子能量的相反数（以带顶为零点，向下为正）：
\begin{equation}
\boxed{E_h=-E(k_1)=\frac{17\hbar^2k_1^2}{6m}\approx5.49\times10^{-21}\,\text{J}\approx0.034\,\text{eV}.}
\end{equation}

\textbf{(6) 空穴带图示}

空穴带是电子能带关于能量轴的镜像。空穴能量在带顶（$k_h=0$）处为零，随 $|k_h|$ 增大而增加（向下为正）。在 $k_h=-k_1$ 处标出空穴位置，能量 $E_h\approx0.034\,\text{eV}$。
\end{derivationbox}

\begin{formulabox}[答案汇总]
\begin{center}
\begin{tabular}{lll}
\toprule
\textbf{物理量} & \textbf{数值} & \textbf{符号/方向} \\
\midrule
(1) 波矢 $k_h$ & $-3.98\times10^8\,\text{m}^{-1}$ & 与 $k_1$ 反向 \\
(2) 有效质量 $m_h^*$ & $1.52\times10^{-31}\,\text{kg}\approx0.17m_0$ & 正值 \\
(3) 速度 $v_h$ & $-2.77\times10^5\,\text{m/s}$ & 与电子速度相同 \\
(4) 动量 $p_h$ & $-4.20\times10^{-26}\,\text{N}\cdot\text{s}$ & $p_h=\hbar k_h=m_h^*v_h$ \\
(5) 能量 $E_h$ & $0.034\,\text{eV}$ & 从带顶向下增加 \\
\bottomrule
\end{tabular}
\end{center}
\end{formulabox}

\begin{insightbox}[物理图像]
\begin{itemize}
    \item 空穴是\textbf{准粒子}：它不是真实的粒子，而是对近满带中大量电子集体行为的等效描述。引入空穴后，只需追踪少量空穴的运动，而无需处理大量电子。
    \item 空穴的``正能量''：$m_h^*>0$、$E_h$ 向下增加、带正电荷 $+e$，使得空穴的运动方程与自由正电荷完全一致。这大大简化了半导体价带的输运理论。
    \item 空穴速度等于缺失电子的速度，这意味着电流方向与电子运动方向相反——这正是半导体中空穴导电的物理本质。
\end{itemize}
\end{insightbox}
